\subsubsection{Các biểu đồ so sánh}

\paragraph{Biểu đồ đường cong học}

\figurename~\ref{fig:regression-learning-curves} thể hiện các đường cong học
của từng mô hình.

\begin{figure}[htbp]
  \centering
  \includegraphics[width=0.9\textwidth]{regression-learning-curves}
  \caption{Biểu đồ đường cong học của các mô hình dựa trên MSE từng kích thước
  của tập huấn luyện}
  \label{fig:regression-learning-curves}
\end{figure}

Ta thấy rằng, hầu hết các mô hình cần khoảng 7000 mẫu huấn luyện để đạt MSE cực
tiểu trên tập huấn luyện và tập xác thực.

\paragraph{Biểu đồ phần dư}

\figurename~\ref{fig:regression-residuals-plots} thể hiện biểu đồ phần dư của
các mô hình trên tập kiểm tra

\begin{figure}[htbp]
  \centering
  \includegraphics[width=0.9\textwidth]{regression-residuals}
  \caption{Biểu đồ phần dư của các mô hình trên tập kiểm tra}
  \label{fig:regression-residuals-plots}
\end{figure}

Ta thấy rằng các mô hình hồi quy tuyến tính cơ bản (OLS, MBGD, WLS) có độ khớp
vừa phải, không quá khớp của sự ngắt ngưỡng.

Ngược lại, các mô hình cơ sở phi tuyến (Polynomial, RBF và Fourier) có việc học
sự ngắt ngưỡng khi mô hình rất ít khi dự đoán trên 5.

Trong khi đó, các mô hình chính quy hoá (Lasso, Ridge, Elastic Net) có sự lệch
khi học.

\paragraph{Biểu đồ giá trị dự đoán-giá trị thực}

\figurename~\ref{fig:regression-predicted-actual} thể hiện tương quan giữa giá
trị dự đoán và giá trị thực của các mô hình trên tập kiểm tra.

\begin{figure}[htbp]
  \centering
  \includegraphics[width=0.9\textwidth]{regression-predicted-actual}
  \caption{Biểu đồ giá trị dự đoán-giá trị thực trên tập kiểm tra}
  \label{fig:regression-predicted-actual}
\end{figure}

Ta thấy rằng các mô hình chính quy hoá cho thấy các điểm phân tán lệch hẳn so
với đường tham chiếu. Điều này cho thấy các mô hình này học chưa khớp khá nặng.

Ngược lại, các mô hình cơ sở phí tuyến có sự khớp cao với cả dữ liệu và việc
ngắt ngưỡng. Điều này có thể ảnh hưởng đến dự đoán thực tế khi không có sự ngắt
ngưỡng này.

Đối với các mô hình hồi quy tuyến tính cơ bản, các điểm phân tán hợp lí hơn. Mô
hình không quá bị ảnh hưởng bởi việc ngắt ngưỡng của dữ liệu. Điều này cho thấy
mô hình khớp tốt với dữ liệu và có khả năng dự đoán thực tế ổn định hơn.
